\chapter{Logica Proposizionale}

Abbiamo visto nel capitolo precedente che una logica è caratterizzata da tre componenti fondamentali: un linguaggio formale, una semantica e un calcolo.
In questo capitolo studieremo in dettaglio la \textit{logica proposizionale}, la più semplice tra le logiche classiche, che costituisce il punto di partenza per comprendere sistemi logici più complessi.
Procederemo sistematicamente, partendo dalla definizione del linguaggio (sintassi), per poi introdurre la semantica e infine studiare i sistemi di calcolo che permettono di derivare conseguenze logiche.

\section{Sintassi della Logica Proposizionale}

Il primo passo nella costruzione di una logica formale è definire il suo linguaggio, specificando quali simboli possiamo usare e come possiamo combinarli per formare espressioni ben formate, detto alfabeto.
In generale, un linguaggio logico è composto da due tipi di simboli: i \textit{simboli non logici}, su cui la logica non fa alcuna assunzione e il cui significato può essere interpretato in modi diversi a seconda del contesto (nel caso della logica proposizionale sono variabili booleane), e i \textit{simboli logici}, i cui significati sono fissati dalla specifica logica.

\begin{definition}{Alfabeto della logica proposizionale}
  L'alfabeto della logica proposizionale è costituito da:
  \begin{itemize}
    \item Un insieme potenzialmente infinito di \textit{simboli non logici}, detti \keyword{proposizioni atomiche}, denotati con \(AP = \set{A,B,C,\ldots}\).
    \item Un insieme di \textit{simboli logici}, che includono i connettivi logici \(\set{\land, \lor, \rightarrow, \neg,\leftrightarrow}\) e le parentesi \(\set{(,)}\) per disambiguare la struttura delle formule.
  \end{itemize}
\end{definition}

Un'\keyword{espressione} è una sequenza finita di simboli del linguaggio.
Tuttavia, non tutte le espressioni sono significative dal punto di vista logico: dobbiamo distinguere quelle che rispettano le regole grammaticali del linguaggio.

\begin{definition}{Formula ben formata (WFF)}
  Una formula ben formata (WFF) è un'espressione che può essere costruita usando una grammatica definita sui simboli del linguaggio.
\end{definition}

La definizione precisa di quali espressioni sono formule ben formate nella logica proposizionale è data ricorsivamente attraverso le seguenti regole di formazione.

\begin{definition}{WFF Proposizionale}
  Una formula è ben formata in logica proposizionale se è costruita secondo le seguenti regole:
\begin{itemize}
  \item \(AP \subseteq WFF\)
  \item Se \(\phi \in WFF\), allora \(\neg \phi \in WFF\)
  \item Se \(\phi,\psi \in WFF\), allora \((\phi \bigodot \psi)\in WFF\), dove \(\bigodot \in \set{\land, \lor, \rightarrow, \leftrightarrow}\) 
  \item Nessuna altra espressione è una formula ben formata
\end{itemize}
  
\end{definition}

Come si può notare, questa definizione induttiva è una grammatica libera da contesto, e di conseguenza, ogni formula ben formata può essere rappresentata come un albero sintattico, dove i nodi interni rappresentano i connettivi logici e le foglie rappresentano le proposizioni atomiche.
Ovviamente, tale albero sintattico cattura la struttura gerarchica della formula, mostrando come essa sia stata costruita a partire dalle sue componenti più semplici.

\begin{example}{Albero sintattico}
  Consideriamo la formula \(\phi = (A \rightarrow (B \land \neg C))\). 
  
  Questa formula è ben formata secondo la definizione ricorsiva: a partire dalle variabili proposizionali \(A\), \(B\), \(C\) (che sono WFF per la prima regola), costruiamo \(\neg C\) (seconda regola), poi \((B \land \neg C)\) (terza regola), e infine \((A \rightarrow (B \land \neg C))\) (terza regola).
  
  L'albero sintattico di \(\phi\) rappresenta questa costruzione gerarchica:
  
  \begin{center}
    \begin{forest}
      [\(\rightarrow\)
        [\(A\)]
        [\(\land\)
          [\(B\)]
          [\(\neg\)
            [\(C\)]
          ]
        ]
      ]
    \end{forest}
  \end{center}
  
  In questo albero:
  \begin{itemize}
    \item Le foglie sono le variabili proposizionali \(A\), \(B\), \(C\)
    \item I nodi interni sono i connettivi logici \(\rightarrow\), \(\land\), \(\neg\)
    \item La struttura dell'albero riflette l'ordine di applicazione delle regole di costruzione delle WFF
  \end{itemize}
\end{example}

Tale costruzione ci permette di utilizzare un potente strumento di ragionamento: l'\textit{induzione strutturale}.
L'induzione strutturale ci permette di dimostrare proprietà delle formule procedendo per casi: il caso base tratta le formule atomiche (variabili proposizionali), mentre i passi induttivi si basano sulla costruzione della formula a partire da sottoformule più semplici.

A differenza dell'induzione matematica semplice sui numeri naturali (dove per dimostrare \(P(n)\) assumiamo solo \(P(n-1)\)), l'induzione strutturale è basata sul principio di \textit{induzione completa} (o forte), dove per dimostrare una proprietà per una formula \(\phi\) assumiamo che valga per \textbf{tutte} le sue sottoformule proprie.
Tale metodo è ben fondato proprio sfruttando l'albero sintattico finito di una formula, essendo l'induzione strutturale equivalente a un'induzione completa su tale albero.

Inoltre, questa definizione induttiva impone forti vincoli sulla forma dei sottoalberi sintattici, che devono essere a loro volta alberi sintattici di formule ben formate.
Tali sottoalberi rappresentano le \keyword{sottoformule} di una formula, il cui insieme è definito come segue.

\begin{definition}{Sottoformule}
  Dato \(\phi \in WFF\), l'\keyword{insieme delle sottoformule} di \(\phi\), denotato con \(Sub(\phi)\), è definito ricorsivamente come segue:
  \[Sub(\phi) = \begin{cases}
    \{\phi\} & ,\phi \in AP\\
    Sub(\phi') \cup \{\neg \phi'\} & ,\phi = \neg \phi'\\
    Sub(\phi') \cup Sub(\psi') \cup \{(\phi' \bigodot \psi')\} & ,\phi = (\phi' \bigodot \psi')
  \end{cases}
  \]
\end{definition}

Possiamo ora osservare una prima applicazione dell'induzione strutturale, dimostrando una proprietà fondamentale delle formule ben formate: il numero di sottoformule di una formula \(\phi\) è al più pari alla lunghezza della formula stessa.

\begin{note}{}
  Tale proprietà evidenzia l'importanza delle regole di formazione delle WFF, differenziandole da generiche successioni di simboli.
  
  Infatti, data una stringa qualsiasi, il numero di sottostringhe è quadratico rispetto alla lunghezza della stringa, mentre per le WFF, grazie alla loro struttura ben definita, il numero di sottoformule è limitato linearmente dalla lunghezza della formula.
\end{note}

\begin{proposition}{}
  Sia \(\phi \in WFF\). Siano inoltre \(\abs{\phi}\) la lunghezza della formula \(\phi\) in numero di simboli (parentesi escluse) e \(\abs{Sub(\phi)}\) la cardinalità dell'insieme delle sottoformule di \(\phi\). Allora
  \[\abs{Sub(\phi)}\leq \abs{\phi}\]
\end{proposition}

\begin{proof}
  Per dimostrare questa proposizione, procediamo per induzione strutturale sulla formula \(\phi\).
  \begin{description}
    \item[\(\phi \in AP\)]: In questo caso, \(\phi\) è una formula atomica, quindi \(Sub(\phi) = \{\phi\}\) e \(\abs{Sub(\phi)} = 1\). Dato che \(\abs{\phi} = 1\) (ogni simbolo ha lunghezza 1), abbiamo
          \[\abs{Sub(\phi)} = 1 \leq 1 = \abs{\phi}\]
    \item[\(\phi = \neg \phi'\)]: In questo caso, per induzione strutturale, assumiamo che \(\abs{Sub(\phi')} \leq \abs{\phi'}\). Dato che \(Sub(\phi) = Sub(\phi') \cup \{\neg \phi'\}\), abbiamo
          \[\abs{Sub(\phi)} = \abs{Sub(\phi')} + 1\]
          Inoltre, \(\abs{\phi} = \abs{\phi'} + 1\) (perché aggiungiamo il simbolo di negazione). Quindi,
          \[\abs{Sub(\phi)} = \abs{Sub(\phi')} + 1 \leq \abs{\phi'} + 1 = \abs{\phi}\]
    \item[\(\phi = (\phi' \bigodot \psi')\)]: Infine, sempre per induzione strutturale, assumiamo che \(\abs{Sub(\phi')} \leq \abs{\phi'}\) e \(\abs{Sub(\psi')} \leq \abs{\psi'}\). Dato che \(Sub(\phi) = Sub(\phi') \cup Sub(\psi') \cup \{(\phi' \bigodot \psi')\}\), abbiamo
          \[\abs{Sub(\phi)} \leq \abs{Sub(\phi')} + \abs{Sub(\psi')} + 1\]
          In questo caso non è detto valga l'uguaglianza, perché è possibile che \(Sub(\phi')\) e \(Sub(\psi')\) non siano disgiunti.
          Inoltre, \(\abs{\phi} = \abs{\phi'} + \abs{\psi'} + 1\) (perché aggiungiamo il connettivo). Quindi partendo da:
          \[\abs{Sub(\phi')} \leq \abs{\phi'}\]
          essendo numeri naturali, possiamo sommare \(\abs{Sub(\psi')}\) a entrambi i membri dell'inequazione senza invertire il verso:
          \[\abs{Sub(\phi')} +\abs{Sub(\psi')} \leq \abs{\phi'} + \abs{Sub(\psi')} \leq \abs{\phi'} + \abs{\psi'}\]
          e dunque
          \[\abs{Sub(\phi)} \leq \abs{Sub(\phi')} + \abs{Sub(\psi')} + 1 \leq \abs{\phi'} + \abs{\psi'} + 1 = \abs{\phi}\]
  \end{description}
\end{proof}

\section{Semantica Proposizionale tramite Modelli}

Definita la sintassi del linguaggio della logica proposizionale, per formalizzare i concetti anticipati nello scorso capitolo è necessario definire formalmente un modo per dare significato alle parole del linguaggio descritto, ovvero la sua \keyword{semantica}.\footnote{La semantica della logica proposizionale formalizza il concetto intuitivo di \q{situazione} in cui una formula è vera o falsa menzionato precedentemente.}

Un modo per definire tale semantica è attraverso il concetto di \keyword{modello}. Un modello, in generale, non è altro che una struttura matematica che interpreta i simboli del linguaggio e stabilisce le condizioni di verità per le formule.
In particolare, per la logica proposizionale, possiamo definire un modello come segue.

\begin{definition}{Modello per la logica proposizionale}
  Un modello \(m\) per la logica proposizionale è un qualsiasi sottoinsieme di \(AP\). In altre parole, \(m \in 2^{AP}\).\footnote{\(2^{AP}\) non è altro che una notazione alternativa per l'insieme di tutti i sott'insiemi di \(AP\), chiamato anche insieme potenza indicato anche con \(\mathcal{P}(AP)\).}
\end{definition}

Essendo la logica proposizionale un linguaggio molto semplice, per darle significato è sufficiente una struttura altrettanto semplice: un insieme di proposizioni atomiche che consideriamo vere.

Essendo \(AP\) potenzialmente infinito, anche i modelli possono essere infiniti. Da un punto di vista matematico non è particolarmente importante, ma modelli infiniti portano a problemi computazionali di rappresentazione e manipolazione.

\begin{definition}{Soddisfacimento di una formula da parte di un modello}
  Definiamo la \keyword{relazione di soddisfazione} \(\models \subseteq 2^{AP} \times PL\) dato un modello \(m \in 2^{AP}\) e una formula \(\phi \in PL\) come segue:
  \(m\models \phi \) se e solo se vale uno dei seguenti casi:
  \begin{enumerate}
    \item \(\phi \in m\) se \(\phi \in AP\).
    \item \(m \not\models \phi'\) se \(\phi = \neg \phi'\).
    \item \(m \models \phi_1\) e \(m \models \phi_2\) se \(\phi = \phi_1 \land \phi_2\).
    \item \(m \models \phi_1\) o \(m \models \phi_2\) se \(\phi = \phi_1 \lor \phi_2\).
    \item Se \(m \models \phi_1\) allora \(m \models \phi_2\) se \(\phi = \phi_1 \rightarrow \phi_2\).
    \item \(m \models \phi_1\) se e solo se \(m \models \phi_2\) se \(\phi = \phi_1 \leftrightarrow \phi_2\).
  \end{enumerate}
\end{definition}

Tale definizione rappresenta una procedura ricorsiva per determinare se una formula è soddisfatta da un modello, partendo dalle proposizioni atomiche e risalendo la struttura sintattica della formula.

Il problema decisionale di verificare se \(m \models \phi\) è detto \keyword{model checking}, e, a patto di risolvere il problema di rappresentazione di modelli potenzialmente infiniti, è facilmente decidibile, poiché si tratta di un'esplorazione in post-ordine dell'albero sintattico di \(\phi\), che ha un numero di nodi pari al numero di sottoformule di \(\phi\), che è al più pari alla lunghezza di \(\phi\).

Avendo finalmente definito la semantica possiamo andare a formalizzare le definizioni informali date all'inizio del corso.
\begin{definition}{Soddisfacibilità}
  Una formula \(\phi\) è \keyword{soddisfacibile} (\(SAT(\phi)\)) se e soltanto se esiste un modello \(m\in 2^{AP}\) tale che \(m \models \phi\).
\end{definition}

Il problema decisionale di verificare se \(SAT(\phi)\) è detto \keyword{soddisfacibilità}, ed è un problema centrale in logica e informatica teorica, poiché molte questioni di interesse possono essere ridotte a problemi di soddisfacibilità.
Rispetto al model checking, il problema di soddisfacibilità è spesso più complesso, poiché richiede l'esplorazione dello spazio dei modelli, che cresce esponenzialmente con il numero di proposizioni atomiche coinvolte nella formula.

\begin{definition}{Validità}
  Una formula \(\phi\) è valida (\(Valid(\phi)\)) se e soltanto se per ogni modello \(m\in 2^{AP}\) vale \(m \models \phi\).
\end{definition}

La validità e la soddisfacibilità sono concetti duali, nel senso che \(SAT(\phi)\iff \neg Valid(\neg \phi)\) e \(Valid(\phi) \iff \not SAT(\neg \phi)\).\footnote{I simboli \(\iff\) e \(\implies\) sono simboli di \textbf{meta-linguaggio}. Tali simboli non appartengono al linguaggio ogetto del discroso (in questo caso la logica proposizionale), ma sono simboli sono utilizzati (con l'interpretazione usuale) per parlare di proprietà del linguaggio stesso.}
Questo è diretta conseguenza del fatto che, data una qualsiasi proprietà \(P\), \(\exists x:P(x) \iff \forall x, \neg P(\neg x)\). Ciò verrà dimostrato formalmente quando tratteremo la logica del primo ordine.

\begin{definition}{Estensione della soddisfacibilità}
  Dato \(\Gamma \subseteq PL\) e \(m \in 2^{AP}\), diremo che \(m\) soddisfa \(\Gamma\) (\(m \models \Gamma\)) se e soltanto se \(m \models \phi\) per ogni \(\phi \in \Gamma\).
\end{definition}

\begin{definition}{Conseguenza logica}
  Dato \(\Gamma \subseteq PL\), una formula \(\phi\) è conseguenza logica di \(\Gamma\) (\(\Gamma \models \phi\)) se e soltanto se per ogni \(m \in 2^{AP}\) che soddisga \(\Gamma\) allora \(m \models \phi\).

  Alternativamente, \(\Gamma \models \phi\) se e soltanto se non esiste \(m \in 2^{AP}\) tale che \(m \models \Gamma\) e \(m \not\models \phi\).
\end{definition}

\begin{note}{}
  Nonostante si usi la stessa notazione \(\models\) per indicare sia la soddisfacibilità di una formula da parte di un modello, sia la conseguenza logica di una formula da parte di un insieme di formule, è importante sottolineare che si tratta di due concetti distinti.
  La soddisfacibilità \(m \models \phi\) è una relazione conenuta in \(2^{AP}\times PL\), mentre la conseguenza logica \(\Gamma \models \phi\) è una relazione contenuta in \(2^{PL}\times PL\).
\end{note}

Dalla definizione di conseguenza logica segue che se \(\emptyset \models \phi\), denotato alternativamente con \(\models \phi\), allora \(\phi\) è valida.
Questo poiché, preso \(\Gamma = \emptyset\), ogni modello soddisfa \(\Gamma\).
Quindi, la relazione di conseguenza logica dall'insime vuoto si riduce alla soddisfacibilità di \(\phi\) da parte di ogni modello, che è proprio la definizione di validità.
In altre parole potendo risolvere il problema decisionale di conseguenza logica, potremmo risolvere anche il problema decisionale di validità.

Inoltre se \(\Gamma\) è un insieme finito di formule vale anche il viceversa, poiché è possibile considereare 
\[(\bigwedge_{\gamma \in \Gamma} \gamma) \rightarrow \phi\]
come una formula singola, la cui validità è equivalente alla conseguenza logica di \(\phi\) da parte di \(\Gamma\).

Ciò, oltre a evidenziare la stretta connessione tra conseguenza logica e validità (e quindi soddisfacibilità per dualità), sottolinea come la relazione \(\models\) sia correlato alla semantica del simbolo di implicazione \(\rightarrow\).

\section{Semantica Proposizionale tramite Assegnamenti}

Definito formalmente un primo concetto di semantica, andiamo a risolvere il problema della sua rappresentabilità.
Per fare ciò, andremo a definire una rappresenzatione più compatta della semantica di una formula, andando a dimostrare come questa sia equivalente alla semantica basata sui modelli, ma con il vantaggio di essere rappresentabile in modo finito.

In primo luogo, per risolvere il problema della potenziale infinitezza dei modelli, possiamo limitare la nostra attenzione solo alle proposizioni atomiche che effettivamente compaiono nella formula di interesse.

\begin{definition}{Proposizioni atomiche di una formula}
  Data \(\phi \in PL\), possiamo definire \(AP(\phi) = Sub(\phi) \cap AP\) l'insieme delle proposizioni atomiche contenute in \(\phi\)
\end{definition}

Essendo che \(\abs{Sub(\phi)}\leq \abs{\phi}\) avremo che \(\abs{AP(\phi)}\leq \abs{\phi}\), che essendo la lunghezza di una parola finita è finito.

\begin{definition}{Assegnamento}
  Dato \(\phi \in PL\), un assegnamento \(\alpha\) per \(\phi\) è una funzione \(\alpha: \Delta_\phi \rightarrow \{0,1\}\), tale che \(\Delta_\phi\) è finito e \(AP(\phi) \subseteq \Delta_\phi \subseteq AP\).

  Chiameremo l'insieme di tali assegnamenti \(Ass(\phi)\).\footnote{Analogamente alla notazione \(2^S\) per indicare l'insieme \(\mathcal{P}(S)\), una notazione molto in uso per indicare un insieme di funzioni \(f: D \to C\) è \(C^D\) essendo che la cardinalità di tale insieme (per insiemi finiti) è proprio \(\abs{C}^{\abs{D}}\). Di conseguenza, una notazione alternativa per \(Ass(\phi)\) è \(\set{0,1}^{\Delta_\phi}\)}
\end{definition}

Tali funzioni di assegnamento ci permettono di assegnare il valore di verità a ciascuna proposizione atomica. Data la natura induttiva della semantica, possiamo definire induttivamente la \keyword{valutazione} del valore di verità di una formula induttivamente partendo da un assegnamento.

\begin{definition}{Valutazione di una formula}
  Siano \(\phi \in PL\) e \(\alpha \in Ass(\phi)\). Definiamo dunque la \keyword{valutazione basata su \(\alpha\)} come una funzione \(val^{\alpha}: PL \to \set{0,1}\) come\footnote{Per essre precisi, \(val\) ha una dipendenza implicita dall'assegnamento, e quindi sarebbe più corretto descriverla come una funzione \(val: Ass(\phi) \to PL \to \set{0,1}\). Nel pratico però, l'assegnamento sarà fissato a priori.}:
  \[val^{\alpha}(\phi) = \begin{cases}
  \alpha(\phi) &, \phi \in AP \\
  1 - val^{\alpha}(\phi') &, \phi = \neg \phi' \\
  \min\set{val^{\alpha}(\phi_1), val^{\alpha}(\phi_2)} &, \phi = \phi_1 \land \phi_2 \\
  \max\set{val^{\alpha}(\phi_1), val^{\alpha}(\phi_2)} &, \phi = \phi_1 \lor \phi_2 \\
  1 & \phi = \phi_1 \rightarrow \phi_2 \text{ e } val^{\alpha}(\phi_1) \leq val^{\alpha}(\phi_2) \\
  0 & \phi = \phi_1 \rightarrow \phi_2 \text{ e } val^{\alpha}(\phi_1) > val^{\alpha}(\phi_2)
  \end{cases}\]
  % \begin{enumerate}
  %   \item Se \(\phi \in AP\), ovvero \(\phi \in \Delta_\phi\), allora \(\alpha(\phi) = val^{\alpha}(\phi)\).
  %   \item Se \(\phi = \neg \phi'\) allora \(val^{\alpha}(\phi) = 1 - val^{\alpha}(\phi')\).
  %   \item Se \(\phi = \phi_1 \land \phi_2\) allora \(val^{\alpha}(\phi) = \min\set{val^{\alpha}(\phi_1), val^{\alpha}(\phi_2)}\).
  %   \item Se \(\phi = \phi_1 \lor \phi_2\) allora \(val^{\alpha}(\phi) = \max\set{val^{\alpha}(\phi_1), val^{\alpha}(\phi_2)}\).
  %   \item Se \(\phi = \phi_1 \rightarrow \phi_2\) allora
  %             \[val^{\alpha}(\phi) = \begin{cases}
  %               1 & val^{\alpha}(\phi_1) \leq val^{\alpha}(\phi_2)\\
  %               0 & \text{ altrimenti}
  %             \end{cases}\]
  % \end{enumerate}
\end{definition}

Varificare se una certa formula assume valore \(0\) o \(1\) rispetto ad un assegnamento compatibile con la formula diventa un problma facilmente decidibile, poiché non si pone più il problema di rappresentabilità di modelli potenzialmente infiniti.
In particolare ora abbiamo una funzione tra domini finiti, delle quali ne esistono \(\abs{\set{0,1}}^{\abs{AP(\phi)}}=2^{\abs{AP(\phi)}}\).

Di consueguenza, il problema di valutazione di un espressione ha complessità \(\BigTheta{\abs{\phi}}\).

Dobbiamo dunque mostrare che questa nuova semantica basata sugli assegnamenti sia equivalente alla semantica basata sui modelli, in altre parole mostrare che a ogni modello corrisponde un assegnamento equivalente e viceversa, e che la relazione di soddisfacimento \(m \models \phi\) è equivalente alla valutazione \(val^{\alpha}(\phi) = 1\).
Per fare ciò, andremo a definire la proprietà di essere una tautologia, che rappresenta la proprietà di una formula di essere valutata a \(1\) a prescindere dall'assegnamento scelto, e mostrare che questa è equivalente alla validità, ovvero alla proprietà di essere soddisfatta da ogni modello.

Facendo ciò, otterremo anche un modo per correlare \(m \models \phi\) e \(val^{\alpha}(\phi) = 1\).

\begin{theorem}{}
  Data \(\phi \in PL\), \(\models \phi \) se e solo se \(Taut(\phi)\)
\end{theorem}

\begin{proof}
  Dalla definizione di validità abbiamo che 
  \[\models \phi \iff \forall m, m \models \phi\]
  e dalla definizione di tautologia si ha che 
  \[Taut(\phi)\iff\forall \alpha val^{\alpha}(\phi) = 1\]

  In altre parole, per avere la tesi, basta mostrare che \(\forall m, m \models \phi \iff \forall \alpha, val^{\alpha}(\phi) =1\).
  Procediamo un verso alla volta:
  \begin{description}
     \item[\(\impliedby\)] %Assumiamo che \(\forall m, m \models \phi\).
     Possiamo notare che \(\forall m,\exists alpha : m\models \phi \iff val^{\alpha}(\phi) = 1\).
      Questo poiché, preso un \(m \in 2^{AP}\) generico posso costruire 
      \[\alpha_{m}(P) = \begin{cases}
        1, & P \in m\\
        0, & \text{ altrimenti}
      \end{cases}\]
    
      Per verificare che \(m\models \phi \iff val^{\alpha_m}(\phi) = 1\) procediamo per induzione strutturale:
      \begin{description}
        \item[\(\phi \in AP\):] In tal caso
          \[m \models \phi \iff \phi \in m \overset{def}{\iff}\alpha_m(\phi) = 1 \iff val^{\alpha}(\phi)\]
        \item[\(\phi = \neg \phi'\):] In tal caso
          \[m \models \phi \iff m \not\models \phi' \overset{ind. hp.}{\iff} val^{\alpha_m}(\phi') = 0\]
          Ma allora dalla definizione di valutazione, \(val^{\alpha_m}(\phi) = 1-val^{\alpha_m}(\phi') = 1-0 = 1\).
        \item[\(\phi = \phi_1 \land \phi_2\):] In tal caso 
          \[m \models \phi \iff m \models \phi_1 \text{ e } m \models \phi_2 \overset{ind. hp.}{\iff} val^{\alpha_m}(\phi_1) = 1 \text{ e } val^{\alpha_m}(\phi_2) = 1\]
          Ma allora dalla definizione di valutazione, \(val^{\alpha_m}(\phi) = \min\set{val^{\alpha_m}(\phi_1), val^{\alpha_m}(\phi_2)} = \min\set{1, 1} = 1\).
        \item[\(\phi = \phi_1 \lor \phi_2\):] In tal caso 
          \[m \models \phi \iff m \models \phi_1 \text{ o } m \models \phi_2 \overset{ind. hp.}{\iff} val^{\alpha_m}(\phi_1) = 1 \text{ o } val^{\alpha_m}(\phi_2) = 1\]
          Ma allora dalla definizione di valutazione, \(val^{\alpha_m}(\phi) = \max\set{val^{\alpha_m}(\phi_1), val^{\alpha_m}(\phi_2)} = \max\set{1, 0} = 1\).
        \item[\(\phi = \phi_1 \rightarrow \phi_2\):] In tal caso 
          \[m \models \phi \iff \text{ se } m \models \phi_1 \text{ allora } m \models \phi_2\]
          Per ipotesi induttiva, se \(m \not\models \phi_1\) allora \(val^{\alpha_m}(\phi_1) = 0\), e quindi a prescindere dal valore di \(val^{\alpha_m}(\phi_2)\) si ha \(val^{\alpha_m}(\phi) = 1\), essendo che \(0 \leq val^{\alpha_m}(\phi_2)\).
          Se invece \(m \models \phi_1\) allora \(val^{\alpha_m}(\phi_1) = 1\), ma sempre per ipotesi induttiva \(m \models \phi_2\) e quindi \(val^{\alpha_m}(\phi_2) = 1\), e dunque \(val^{\alpha_m}(\phi) = 1\), essendo che \(1 \leq 1\).
      \end{description}

      Dunque, per ogni modello \(m\) esiste un assegnamento \(\alpha_m\) tale che la sua valutazione di \(\phi\) è \(1\) solo se \(m\) soddisfa \(\phi\).
      Andiamo dunque a dimostrare l'implicazione da destra a sinistraper contrapposizione, ovvero mostrando che, assumendo \(\not\models \phi\) segue che \(\neg Taut(\phi)\).
      Ma allora \(\exists m, m \not\models \phi \iff m \models \neg\phi\).
      
      Dalla costruzione precedente però sappiamo che, fissato questo \(m\), esiste \(\alpha_m : m \models \neg \phi \iff val^{\alpha_m}(\neg\phi) = 1\).
      Ma allora \(val^{\alpha_m}(\neg \phi) = 1\) implica \(val^{\alpha_m}(\phi) = 0\), e dunque \(\neg Taut(\phi)\).

  \item[\(\impliedby\)] anche qui per contrapposizione assumo \(\neg Taut(\phi)\) e dimostro \(\not\models\phi\).
    Per fare questo ci serve mostrare che \(\forall \alpha \exists m_{\alpha}: val^{\alpha}(\phi) = 1 \iff m\models\phi\).
    Ma allora basta definire \(m_\alpha = \set{P \in AP(\phi)}[\alpha(P) = 1]\).
    
    In maniera analoga al caso precedente, possiamo dimostrare che \(val^{\alpha}(\phi) = 1 \iff m_\alpha \models \phi\) per induzione strutturale su \(\phi\).

    A questo punto però, assumendo \(\neg Taut(\phi)\) abbiamo che \(\exists \alpha, val^{\alpha}(\phi) = 0\), e dunque \(val^{\alpha}(\neg \phi) = 1\), e quindi \(m_\alpha \models \neg \phi\), e dunque \(m_\alpha \not\models \phi\), e quindi \(\not\models \phi\).
  \end{description}
\end{proof}

Questo risultato ha come sottoprodotto il fatto che \(Valid\) è computabile, poiché è equivalente a \(Taut\), che è computabile, essendo su insiemi finiti.
Il costo computazionale di \(Valid\) è però esponenziale, poiché il \(\abs{Ass(\phi)} = 2^{\abs{\phi}}\), e dunque \(Valid\) è \(\BigO{\abs{\phi}2^{\abs{\phi}}}\leq 2^{\BigTheta{\abs{\phi}}}\).

\begin{example}{Tautologie}
  Vediamo alcune tautologie notevoli.
  \begin{enumerate}
    \item \(P \lor \neg P\) (legge del terzo escluso).
    \item \(\neg (P \land \neg P)\) (principio di non contraddizione \(\bot\))
    \item \(P \lor Q \leftrightarrow \neg(\neg P \land \neg Q)\) (legge di De Morgan)
    \item \(P\land Q \leftrightarrow \neg(\neg P \lor \neg Q)\) (legge di De Morgan)
    \item \(\neg \neg P \iff P\)
    \item \(P \land(Q \lor R) \leftrightarrow (P \land Q) \lor (P \land R)\) (distributività)
    \item \((P \lor(Q\land R)) \leftrightarrow (P \land (Q \lor R))\) (distributività)
    \item \((P \rightarrow Q) \leftrightarrow (\neg P \lor Q)\) (implicazione materiale)
  \end{enumerate}
\end{example}

\subsection{Sostituzioni}

\todo{Aggiungi piccola prefazione per contestualizzare. Probabilmente richiede prossima lezione}

\begin{definition}{Sostituzione}
  Definiamo una \keyword{sostituzione} come una funzione \(\sigma: AP \to PL\).
\end{definition}

Tale funzione ci permette di sostituire ogni proposizione atomica con una formula arbitraria, e dunque di costruire nuove formule a partire da formule già esistenti.

\begin{definition}{Applicazione di una sostituzione a una formula}
  Dato \(\phi \in PL\) e una sostituzione \(\sigma\), definiamo l'applicazione di \(\sigma\) a \(\phi\), denotata con \(\phi\sigma\), come segue:
  \begin{itemize}
   \item se \(\phi \in AP\) allora \(\phi\sigma = \sigma(\phi)\)
   \item se \(\phi = \neg \phi'\) allora \(\phi\sigma = (\neg \phi ')\sigma = \neg (\phi'\sigma)\)
   \item se \(\phi = \phi_1 \bigcirc \phi_2\) allora \(\phi\sigma = (\phi_1 \bigcirc \phi_2)\sigma = (\phi_1\sigma) \bigcirc (\phi_2\sigma)\)
  \end{itemize}
\end{definition}

\begin{note}{}
  È importante sottolineare che la sostituzione è da intendersi come applicata simultaneamente a tutte le proposizioni atomiche. Se una proposizione atomica \(P\) compare sia come simbolo da sostituire che come simbolo di una formula sostituita, la formula sostituita a \(P\) non viene ulteriormente sostituita anche internamente alle formule sostituite contenenti \(P\).
\end{note}

\begin{example}{Sostituzioni e loro applicazione}  
  Consideriamo \(\sigma = \set{A \leftarrow (A\lor B), B \leftarrow \neg A, C \leftarrow C}\) e \(\phi = (A \land B) \lor C\).
  
  Allora, \(\phi\sigma\) sarà \[((A\lor B) \land \neg A) \lor C\]
\end{example}
